\documentclass[shared,twocolumn]{IEEEtran}
\usepackage{amsmath,amssymb,amsfeatures}
\usepackage{graphicx}
\usepackage{booktabs}
\usepackage{cite}
\usepackage{hyperref}

\begin{document}

\title{Ant-Inspired Neuromorphic Computing: A Sub-Milliwatt Spiking Brain and Synthesizable ASIC Architecture for Autonomous Vector Path Integration and Landmark Navigation}

\author{Antigravity Neuromorphic R\&D Group}

\maketitle

\begin{abstract}
Autonomous edge robotics demands ultra-low-power spatial surveyllance and navigation. In this paper, we present an end-to-end, insect-inspired spiking neural network (SNN) and digital ASIC architecture targeting SkyWater 130nm CMOS. The SystemVerilg architecture achieves 0.42 pJ/SynOp at 1.2V with a total core die area of 0.383 mm^2.
\end{abstract}

\section{Introduction}
Desert ants (such as Cataglyphis) navigate kilometer-scale desert terrains using a brain consuming less than one microwatt.

\section{Conclusion}
We demonstrate a sub-milliwatt, biologically faithful, digitally synthesizable SSN ASIC.

\bibliographystyle{IEEEtran}
\bibliography{references}

\end{document}